\begin{flushright}
\huge{\textbf{Abstract}}
\end{flushright}
\addcontentsline{toc}{chapter}{Abstract}

%%%%%%%%%%%%%%%%%%%%%%%%%%%%%%%%%%%%%%%%%%%%%%%%%%%%%%%%%%%%%%%%%%%%%%%%%%%%%%%%%%%%



In the grand academic tradition of IIT Jodhpur, where minds are honed and sleep cycles are forgotten, this thesis stands as a monumental testament to the sleepless nights, countless simulations, failed experiments, and cups of tea that fueled its creation. Under the wise and occasionally mysterious guidance of my supervisor(s), I have ventured into the abyss of research, hoping to emerge with something more than just existential dread and a permanent squint.

This document represents the culmination of the many ups and downs of my Ph.D. journey—mostly downs if we're being honest. Yet, within these pages lies the essence of technical brilliance (I hope), scientific inquiry (mostly), and a few well-timed epiphanies at 3 AM. It is more than a mere collection of chapters; it is a chronicle of academic survival.  It represents my attempt to synthesize years of work, and countless cups of tea into a narrative that is both intellectually rigorous and mercifully brief (for the reader's sake).

Here, you will find the detailed academic tale of occasional existential crisis, which is fondly referred to as "research"- as they say. If I had explained it to my past self, it would have sounded like an obscure sci-fi plot. Through this thesis, I aim to contribute to the vast ocean of knowledge—by at least a tiny drop—and to prove that despite my shaky relationship with proper sleep, my brain can still churn out something helpful and intelligent. Think of it as the ultimate story of your academic adventure—a tale of triumph over tricky hypotheses and of victory against vicious reviewers.

In the end, if this thesis helps to impart even a small bit of education in professional technical communication, then perhaps it wasn’t all just a caffeine-induced hallucination after all.

